\subsection{The relative singularity category retained by Swan lattices}

\begin{lemma}[Proved: the exact relative matrix-factorization model for dyadic Swan lattices]
\label{C:relative-Swan-MF-bridge}
Let $R$ be a commutative noetherian local ring of finite
Krull dimension with $2$ in its maximal ideal.
Put
\[
 B=R[t]/(t^2-1),\qquad Q=R[U],\qquad q=U(U+2),
 \qquad B=Q/(q)\quad(t=U+1).
\]
Let $\mathcal E$ be the exact category of finitely
generated $B$-modules that are free over $R$, with
all short exact sequences between such modules.
It is Frobenius, with projective-injective objects
the finite free $B$-modules.
There are equivalences of triangulated categories
\[
 \underline{\mathcal E}
 \simeq
 \ker\bigl(D_{\rm sg}(B)\longrightarrow D_{\rm sg}(R)\bigr)
 \simeq
 D_{\rm sg}^{\rm rel}(Q\twoheadrightarrow B)
 \simeq [\operatorname{MF}(Q,q)].                    \tag{C.1.RM1}
\]
The kernel in (C.1.RM1) is the full subcategory of
objects whose restriction to $R$ is zero in
$D_{\rm sg}(R)$.
The relative category is the quotient of the bounded
derived category objects that are perfect over $Q$
by $\operatorname{Perf}(B)$, as in
\cite[Definition 2.6]{C-BurkeWalker2012}.
No idempotent completion of these categories is
silently inserted.
In particular,
\[
 R_R(C_2)/\mathbb Z[B]
            \simeq K_0([\operatorname{MF}(R[U],U(U+2))]).
                                                               \tag{C.1.RM2}
\]
For a lattice with involution matrix $T$, its
matrix factorization in this model is
\[
          ((U+1)I-T,\ (U+1)I+T).                    \tag{C.1.RM3}
\]
These are identifications with a relative category;
they do not identify its Grothendieck group with
the group of the ambient singularity category.
Here matrix factorizations use finite projective
$Q$-modules, which are not assumed free.
\end{lemma}
\begin{proof}
Because $R$ is local, every finite projective
$R$-module is free.  Every exact sequence in
$\mathcal E$ therefore splits over $R$.
For $M\in\mathcal E$, its induced module
$B\otimes_R M$ is finite free over $B$.
The action map onto $M$ is an admissible
surjection, and its kernel is the sign twist
$sM$:
\[
 0\longrightarrow sM\longrightarrow B\otimes_R M
       \longrightarrow M\longrightarrow0.           \tag{C.1.RM4}
\]
In coordinates $B\otimes_R M=M\oplus tM$, the
surjection sends $(a,b)$ to $a+Tb$ and the
kernel is parametrized by $b\mapsto(-Tb,b)$.
The action of $t$ on that parameter is $-T$.
Applying the same sequence to $sM$ supplies
an admissible injection of $M$ into an induced
module with cokernel $sM$.

The algebra $B$ is a symmetric Frobenius
$R$-algebra: the coefficient of $1$ in a
product gives a perfect $R$-bilinear pairing
$B\times B\to R$ and hence
$B\simeq\operatorname{Hom}_R(B,R)$ as
$B$-bimodules.
Consequently induced and coinduced modules
agree.  By their adjunctions they are both
projective and injective for $R$-split
sequences.  Any projective or injective
object in $\mathcal E$ is a direct summand
of an induced object, by the preceding
surjection or injection.
It is therefore projective over $B$.
Conversely every finite projective $B$-module
has both relative properties.
The finite algebra $B$ is local: its
quotient by the maximal ideal of $R$ is
$\kappa[t]/((t-1)^2)$, which has one maximal
ideal, and every maximal ideal of $B$
lies over the maximal ideal of $R$.
Thus its finite projectives are free.
This proves the Frobenius assertion.
It also shows that suspension and syzygy
on $\underline{\mathcal E}$ both identify
with sign twist.

The inclusion into $D_{\rm sg}(B)$ is fully
faithful.  Here this can be checked using
the explicit complete resolution
\[
 \cdots\xrightarrow{tI+T}B^d
 \xrightarrow{tI-T}B^d
 \xrightarrow{tI+T}B^d
 \xrightarrow{tI-T}B^d\longrightarrow M\longrightarrow0.
                                                               \tag{C.1.RM5}
\]
For instance, writing an element of $B^d$
as $a+tb$, the equation $(tI-T)(a+tb)=0$
is $b=Ta$.  These elements are exactly
the image of $tI+T$.
The other parity follows in the same way.
The $B$-dual complex is exact by the same
calculation for $T^{\mathsf t}$.
Thus (C.1.RM5) is totally acyclic.
The complete-resolution computation of
singularity-category morphisms
\cite[Lemma B.6]{C-BurkeWalker2012}
identifies its degree-zero Hom cohomology
with the $B$-linear maps modulo those
factoring through finite projective
$B$-modules.  This is precisely the Hom
group in $\underline{\mathcal E}$.
The usual stable triangles come from
the admissible exact sequences, so
the inclusion is triangulated.

Restriction sends a perfect $B$-complex
to a perfect $R$-complex because $B$
is finite free over $R$.
It therefore defines the functor in
(C.1.RM1), and every object of
$\underline{\mathcal E}$ belongs to
its kernel.
Conversely let $X\in D^b(B)$ become
perfect over $R$.
Truncate a bounded-above resolution
of $X$ by finite free $B$-modules
sufficiently far below its cohomology.
In $D_{\rm sg}(B)$ this replaces $X$
by a shift of a finite module $N$.
Restriction of the truncation triangle
shows that $N$ has finite projective
dimension over $R$.
Taking further $B$-syzygies gives an
$R$-projective module: each term of the
$B$-free resolution is $R$-free, so
dimension shifting lowers the finite
$R$-projective dimension to zero.
This syzygy is in $\mathcal E$.
Its shift is again in the image of
$\underline{\mathcal E}$, since sign
twist realizes every suspension.
Thus the essential image is exactly
the displayed kernel, without adding
direct summands to obtain density.

It remains to identify the relative
category correctly.  For a complex
of $B$-modules with bounded finite
cohomology, being perfect over $Q$
is equivalent to being perfect over
$R$.
One implication follows because $Q$
is free over $R$: a finite perfect
$Q$-resolution gives finite $R$-Tor
amplitude, and bounded finite
$R$-cohomology then gives an
$R$-perfect complex.
For the converse use the exact
polynomial-action sequence
\[
 0\longrightarrow Q\otimes_R X
  \xrightarrow{U\otimes1-1\otimes U}
  Q\otimes_R X\longrightarrow X\longrightarrow0.
                                                               \tag{C.1.RM6}
\]
It is termwise exact for complexes;
injectivity follows by comparing the
highest power of the polynomial
variable, and the final map evaluates
that variable at the given action.
If $X$ is perfect over $R$, the first
two terms are perfect over $Q$,
so its cone $X$ is perfect over $Q$.
This proves the equivalence.

Since the monic polynomial $q$ is a
nonzerodivisor on $Q$, the immersion
$\operatorname{Spec}B\hookrightarrow
\operatorname{Spec}Q$ has finite flat
dimension.  Definition 2.6 of
\cite{C-BurkeWalker2012} therefore
identifies its relative singularity
category with precisely the full
kernel just established.
Their Theorem 2.7 applies to
$\operatorname{Spec}Q$, projective
over itself with the trivial very
ample line bundle, and regular
section $q$.
The assumptions that $Q$ is
noetherian and finite dimensional
hold; regularity of $Q$ or a
coefficient field is not required.
It gives the final equivalence
with the homotopy category of
finite-projective matrix
factorizations.
On this affine scheme the source's
quotient by locally contractible
factorizations is the ordinary
homotopy category: a locally
null-homotopic identity is globally
null-homotopic, since global sections
are exact on the quasi-coherent
mapping complex.
Formula (C.1.RM3) follows by direct
multiplication, and its cokernel
is exactly the original lattice.

Finally the Grothendieck group
of a Frobenius stable category
is the Grothendieck group of
its exact category modulo the
classes of projective-injectives.
Indeed those objects become zero,
and the defining triangles are
exactly the triangles obtained
from admissible short exact
sequences and their isomorphs.
Here $K_0(\mathcal E)=R_R(C_2)$
and every projective-injective
is finite free over the local
ring $B$.
This proves (C.1.RM2).
\end{proof}

\noindent\emph{Remaining pencil problem.}
For the node, the regular ambient
ring $S_0[U]$ and the regular
sequence $(g,U(U+2))$ meet the
nonequicharacteristic hypotheses
of \cite[Theorem 1 and Appendix A]
{C-BurkeWalker2012}.
The resulting projective
hypersurface equivalence describes
the ambient $D_{\rm sg}(B)$.
The relevant category remains
its full relative subcategory
in (C.1.RM1).
Neither a restriction of that
equivalence to specified pencil
supports nor injectivity on
$K_0$ of such a subcategory
has been proved here.
An ordinary $G_0$ or ambient
singularity-$K_0$ computation
does not resolve that remaining
relative computation.
